\section{The two differential operators involving the terminal matrix}
\label{rc-section}

We next construct the remaining two polynomial differential
operators in the four-operator argument.  The calculation
extends Section~\ref{rt-section} by retaining a block of
terminal indices.  A general bordered-minor identity gives
both the ordinary recurrence and its boundary term.
All signs and normalizations below refer to the actual
terminal operators of the source \cite{alpo-uniform-source}.

\subsection{The terminal entries and a bordered-minor identity}

Use the notation of Section~\ref{rt-section}, with the
ordered basis
\[
 e_1,\ldots,e_r,\ w_0,\ldots,w_k,\
                         e_r^\vee,\ldots,e_1^\vee .
\]
Here $h(w_i,w_j)=\epsilon_i\delta_{ij}$, where
$\epsilon_0=\cdots=\epsilon_{k-1}=1$ and
$\epsilon_k=-1$.  Put $d=k+1$.
For $m\in U(W)$ set
\[
 a_j=(m^{-1}z)_j,\qquad b_i=h(w_i,m^{-1}z),
 \qquad a'_j=z'_j,\qquad b'_i=h(w_i,z').
\]
These four vectors have exactly the coefficient conventions
in the terminal matrix.  Let $Q^W$ denote the quantum
minors made from right derivatives on $U(W)$.
Write $\widehat i$ for the increasing list
$\{0,\ldots,k\}\setminus\{i\}$, and similarly
$\widehat{i,i'}$ for two omissions.  Set $v=Y+k/2$.
The three entries needed here are
\begin{align}
 t_{21}(Y)&=
 -\sum_{i,j}(-1)^{i+j}b'_i a_j
                    Q^W_{\widehat i,\widehat j}(v),
                                                     \label{rc-t21}\\
 t_{12}(Y)&=
 -\sum_{i,j}(-1)^{i+j}b_i a'_j
                    Q^W_{\widehat i,\widehat j}(v-1),
                                                     \label{rc-t12}\\
 t_{22}(Y)&=
 -\sum_{\substack{i<i'\\j<j'}}(-1)^{i+i'+j+j'}
     (b_i b'_{i'}-b_{i'}b'_i)
     (a_j a'_{j'}-a_{j'}a'_j)
                 Q^W_{\widehat{i,i'},\widehat{j,j'}}(v-1).
                                                     \label{rc-t22}
\end{align}
Every scalar coefficient in these formulas multiplies
the result of the indicated differential operator.
It is not differentiated by that operator.
The last entry is the one already identified with a
positive compact Capelli contraction in
Proposition~\ref{rp-finite-terminal}.

We record two algebraic forms of the same determinant
calculation.  Lists $I=(i_1,\ldots,i_q)$ and
$J=(j_1,\ldots,j_q)$ in the following lemma have distinct
entries, none equal to a distinguished index $e$.
They need not be equal to one another.
All indices in signs inside this lemma start at one.

\begin{lemma}\label{rc-border}
Let $E_{ab}$ satisfy the general linear Lie algebra
relations, and let $\Lambda$ be a complex linear map
from their enveloping algebra to a complex vector space.

Suppose that for every $X$ and all $i\in I,j\in J$,
\[
 \Lambda(E_{ie}X)=\ell_i\Lambda(X),\qquad
 \Lambda(XE_{ej})=r_j\Lambda(X).
\]
Then
\begin{align}
 \Lambda\bigl(Q_{(e,I),(e,J)}(u)\bigr)
 &=\Lambda\bigl((u+E_{ee})Q_{I,J}(u-1)\bigr)\nonumber\\
 &\quad-\sum_{\alpha,\beta=1}^{q}
     (-1)^{\alpha+\beta}\ell_{i_\alpha}r_{j_\beta}
       \Lambda\bigl(Q_{I\setminus i_\alpha,J\setminus j_\beta}
                                                   (u-2)\bigr).
                                                     \label{rc-border-first}
\end{align}
If instead
\[
 \Lambda(E_{ej}X)=\ell_j\Lambda(X),\qquad
 \Lambda(XE_{ie})=r_i\Lambda(X),
\]
then
\begin{align}
 \Lambda\bigl(Q_{(I,e),(J,e)}(u)\bigr)
 &=\Lambda\bigl(Q_{I,J}(u)(u-q+E_{ee})\bigr)\nonumber\\
 &\quad-\sum_{\alpha,\beta=1}^{q}
     (-1)^{\alpha+\beta}r_{i_\alpha}\ell_{j_\beta}
       \Lambda\bigl(Q_{I\setminus i_\alpha,J\setminus j_\beta}
                                                     (u)\bigr).
                                                     \label{rc-border-last}
\end{align}
The formulas include $q=0$, with empty sums.
\end{lemma}

\begin{proof}
Besides its column expression, a quantum minor has
the reverse row expression
\begin{equation}\label{rc-row-form}
 Q_{I,J}(u)=
 \sum_{\tau\in S_q}\operatorname{sgn}(\tau)
 \prod_{\alpha=q}^{1}
 \bigl((u-\alpha+1)\delta_{i_\alpha,j_{\tau(\alpha)}}
                                  +E_{i_\alpha,j_{\tau(\alpha)}}\bigr).
\end{equation}
This follows from the anti-involution
$E_{ab}^*=E_{ba}$ and the adjoint identity in
Theorem~\ref{rp-capelli}: apply that identity to
$Q_{J,I}(\overline u)$ and reverse each product.
It is an identity in the enveloping algebra and
does not require a Hilbert-space realization.

Expand the determinant on the left of
\eqref{rc-border-first} in its first column.
The first term is
$(u+E_{ee})Q_{I,J}(u-1)$.
For row $i_\alpha$ the sign is $(-1)^\alpha$
and the first factor is $E_{i_\alpha,e}$.
Remove it using the left rule for $\Lambda$.
In the remaining minor the first row is $e$.
Expanding \eqref{rc-row-form} in that first row
places $E_{e,j_\beta}$ at the far right and
leaves the minor with parameter $u-2$.
Its sign is $(-1)^{\beta-1}$.
Remove the rightmost factor by the right rule.
The combined sign is
$(-1)^{\alpha+\beta-1}=-(-1)^{\alpha+\beta}$,
which proves \eqref{rc-border-first}.

For \eqref{rc-border-last}, expand in the last
column.  The diagonal term is
$Q_{I,J}(u)(u-q+E_{ee})$.
For row $i_\alpha$ the remaining minor has last
row $e$ and is followed by $E_{i_\alpha,e}$.
The column sign is $(-1)^{\alpha+q+1}$.
Remove the final factor by the right rule.
Expansion of \eqref{rc-row-form} in the last
row puts $E_{e,j_\beta}$ at the far left,
with sign $(-1)^{q+\beta}$; the minor left
over has parameter $u$.
The product of the two signs is again
$-(-1)^{\alpha+\beta}$.
This proves the second identity.
\end{proof}

\subsection{Mixed roots and the boundary contraction}

Let $F$ satisfy the smoothness and two character
laws \eqref{rt-whittaker}.  At a fixed cell point
$g=b(c)m$, evaluation of a right differential word
applied to $F$ gives a map $\Lambda$ as in
Lemma~\ref{rc-border}.
In addition to Lemma~\ref{rt-root-rules}, the
character laws give, for $1\leq l\leq r$,
\begin{align}
 E_{e_l,w_j}F&=\delta_{l,r}a'_jF,&
 E_{w_i,e_l^\vee}F&=\delta_{l,r}b'_iF,
                                                     \label{rc-mixed-right}\\
 (E_{w_i,e_l}XF)(b(c)m)
     &=\delta_{l,r}c_l^{-1}b_i(XF)(b(c)m),&
 (E_{e_l^\vee,w_j}XF)(b(c)m)
     &=\delta_{l,r}\bar c_l^{-1}a_j(XF)(b(c)m)
                                                     \label{rc-mixed-left}
\end{align}
for every right differential word $X$.

Here is a direct verification of the normalizations.
For the root $e_r w_j^*$ the real unitary root
matrix is
$t e_r w_j^*-\bar t\,\epsilon_jw_j(e_r^\vee)^*$.
Its contribution to $\ell_{z'}$ is $t z'_j$,
so its complex character derivative is $a'_j$.
For $w_i(e_r^\vee)^*$ the contribution is
$-\bar t\,\epsilon_i z'_i$, whose character
derivative is $\epsilon_i\overline{z'_i}=b'_i$.
Earlier isotropic indices give zero.
At a cell point conjugation gives
\[
 \operatorname{Ad}(b(c)m)(w_i e_l^*)
      =c_l^{-1}(mw_i)(e_l^\vee)^*,\qquad
 \operatorname{Ad}(b(c)m)(e_l^\vee w_j^*)
      =\bar c_l^{-1}e_l(w_j^*m^{-1}).
\]
The left character derivatives are respectively
$c_l^{-1}h(mw_i,z)=c_l^{-1}b_i$
and $\bar c_l^{-1}(m^{-1}z)_j$ when $l=r$.
Right differentiation preserves left equivariance,
which proves \eqref{rc-mixed-left} with the
arbitrary word $X$.

There is one sign in the boundary term which is
useful to make explicit.  For a fixed omitted
index $s$, write
$p_s(i)=i-\mathbf1_{i>s}$ for the position of
$i$ in the increasing list $\widehat s$,
starting at zero.  Grouping a double sum by
unordered row and column pairs gives
\begin{align}
 &\sum_{s,t}(-1)^{s+t}b'_s a_t
       \sum_{\substack{i\ne s\\j\ne t}}
        (-1)^{p_s(i)+p_t(j)}b_i a'_j
            Q^W_{\widehat{s,i},\widehat{t,j}}(v-1)
       =t_{22}(Y),                                  \label{rc-boundary-sign}\\
 &\sum_{s,t}(-1)^{s+t}b_s a'_t
       \sum_{\substack{i\ne s\\j\ne t}}
        (-1)^{p_s(i)+p_t(j)}b'_i a_j
            Q^W_{\widehat{s,i},\widehat{t,j}}(v-1)
       =t_{22}(Y).                                  \label{rc-boundary-sign-dual}
\end{align}
The lists with two omissions are always put in
increasing order.
For \eqref{rc-boundary-sign}, the two row orders
give $b'_s b_i-b'_i b_s$, which is the negative
of the row factor in \eqref{rc-t22}.
The two column orders give its column factor.
For \eqref{rc-boundary-sign-dual} the row
factor has the correct order and the column
factor changes sign.  In each case the minus
sign is exactly the one in \eqref{rc-t22}.
For $s<i,t<j$, the position corrections
subtract two from $s+i+t+j$ and hence do not
change its parity.  These observations verify
all four terms for every pair and prove
both identities.

\subsection{The ordinary and conjugate recurrences}

For $0\leq l\leq r$, let
\[
 I_{l,s}=(e_{l+1},\ldots,e_r,
                             w_0,\ldots,\widehat w_s,\ldots,w_k),
\]
and
\[
 J_{l,s}=(w_0,\ldots,\widehat w_s,\ldots,w_k,
                             e_r^\vee,\ldots,e_{l+1}^\vee).
\]
Let $\gamma$ and $\Phi$ be as in \eqref{rt-gauge}.
Define
\begin{align}
 F_l(c,m)
  &=\gamma(c)\sum_{s,t}(-1)^{s+t}b'_s a_t\,
       \bigl(Q_{I_{l,s},I_{l,t}}(u_l)F\bigr)(b(c)m),
       &u_l&=Y+r-l+\frac k2,                       \label{rc-F}\\
 F_l^-(c,m)
  &=\gamma(c)\sum_{s,t}(-1)^{s+t}b_s a'_t\,
       \bigl(Q_{J_{l,s},J_{l,t}}(v-1)F\bigr)(b(c)m).
                                                     \label{rc-Fminus}
\end{align}
The coefficients $a,b$ here are evaluated at $m$.
They are independent of $c$.

\begin{proposition}\label{rc-recurrences}
The terminal values are
\[
                F_r=-t_{21}(Y)\Phi,\qquad
                F_r^-=-t_{12}(Y)\Phi.
\]
For $r\geq1$ the boundary formulas are
\begin{align}
 F_{r-1}&=(Y+p_r)F_r-c_r^{-1}t_{22}(Y)\Phi,
                                                     \label{rc-first-boundary}\\
 F_{r-1}^-&=(Y-\bar p_r)F_r^-
                              -\bar c_r^{-1}t_{22}(Y)\Phi .
                                                     \label{rc-last-boundary}
\end{align}
For $0\leq l\leq r-2$,
\begin{align}
 F_l&=(Y+p_{l+1})F_{l+1}
                       -\frac{c_{l+2}}{c_{l+1}}F_{l+2},
                                                     \label{rc-first-bulk}\\
 F_l^-&=(Y-\bar p_{l+1})F_{l+1}^-
                       -\frac{\bar c_{l+2}}{\bar c_{l+1}}F_{l+2}^- .
                                                     \label{rc-last-bulk}
\end{align}
\end{proposition}

\begin{proof}
Right derivatives with indices in $W$ act on
$F(b(c)m)$ as the corresponding right derivatives
in $m$.  The scalar $\gamma$ is independent of
$m$.  The definitions \eqref{rc-t21} and
\eqref{rc-t12} therefore give the terminal values.

For \eqref{rc-first-bulk}, apply
\eqref{rc-border-first} with distinguished
index $e_{l+1}$, before inserting the outside
coefficients in \eqref{rc-F}.
All left and right root constants vanish except
for the next index $e_{l+2}$.  The product of
its two constants is $c_{l+2}/c_{l+1}$.
The remaining parameters are $u_l-1=u_{l+1}$
and $u_l-2=u_{l+2}$.
The diagonal right derivative becomes $p_{l+1}$
on the cell.  Conjugating it by $\gamma$ changes
$u_l+p_{l+1}$ into $Y+p_{l+1}$, since
$p_{l+1}\log\gamma=r-l+k/2$.
The outside coefficients are independent of this
coordinate, so the derivative does not create
an additional term.

For $l=r-1$, the remaining indices are the
$W$-indices with one omission.
Equations \eqref{rc-mixed-right} and
\eqref{rc-mixed-left} give constants
$c_r^{-1}b_i$ and $a'_j$.
The contracted double sum is exactly
$c_r^{-1}t_{22}(Y)$ by
\eqref{rc-boundary-sign}.  The same scalar
conjugation proves \eqref{rc-first-boundary}.

For the conjugate recurrence apply
\eqref{rc-border-last}, with distinguished
index $e_{l+1}^\vee$ appended to the two lists.
Its diagonal generator commutes with every generator
in the remaining minor, since that minor has no row
or column with the distinguished index.  We may therefore
move the diagonal factor to the left before restricting
to the cell.
The right derivative in its diagonal is
$-\bar p_{l+1}$.
If the original size of the minor is
$r-l+k$, its diagonal scalar is
\[
       (v-1)-(r-l+k-1)=Y-r+l-\frac k2 .
\]
Conjugating $-\bar p_{l+1}$ by $\gamma$
adds $r-l+k/2$, leaving $Y-\bar p_{l+1}$.
All parameters of the smaller minors remain
$v-1$, exactly as in \eqref{rc-Fminus}.

Away from the boundary the only surviving
pair of root constants is the adjacent dual
pair, with product
$\bar c_{l+2}/\bar c_{l+1}$.
At the boundary their values are
$b'_i$ and $\bar c_r^{-1}a_j$ by
\eqref{rc-mixed-right} and
\eqref{rc-mixed-left}.
Identity \eqref{rc-boundary-sign-dual}
identifies the double contraction with
$\bar c_r^{-1}t_{22}(Y)$.
This proves the conjugate formulas.
\end{proof}

\begin{corollary}\label{rc-chain}
With the transfer columns of
Theorem~\ref{rt-actual-lifts},
\begin{align}
 F_0&=-t_{21}(Y)A_r(Y)\Phi
                              -t_{22}(Y)B_r(Y)\Phi,\label{rc-chain-plus}\\
 F_0^-&=-t_{12}(Y)\overline A_r(Y)\Phi
                              -t_{22}(Y)\overline B_r(Y)\Phi .
                                                     \label{rc-chain-minus}
\end{align}
\end{corollary}

\begin{proof}
All terminal operators act only on $m$ and commute
with coordinate operators in $c$.
Expand the recurrence \eqref{rc-first-bulk}
over the path of vertices $1,\ldots,r$.
The coefficient of $F_r$ is the sum over
disjoint adjacent edges, with diagonal factors
$Y+p_i$ at uncovered vertices and edge
factors $-c_{i+1}/c_i$.
As in the proof of Theorem~\ref{rt-actual-lifts},
this is $A_r(Y)$.
The coefficient of the other boundary term is
the same sum on $1,\ldots,r-1$, multiplied by
$c_r^{-1}$; it equals
$c_r^{-1}A_{r-1}(Y)=B_r(Y)$.
Substituting the terminal values proves
\eqref{rc-chain-plus}.
The same expansion, with $p_i$ replaced by
$-\bar p_i$ and $c_i$ by $\bar c_i$, gives
\eqref{rc-chain-minus}.
For $r=0$ the assertions reduce to the terminal
identities and $B_0=\overline B_0=0$.
\end{proof}

\subsection{Polynomial coefficients on the full group}

Let $\operatorname{vol}$ be the alternating
form taking value one on the displayed ordered
basis of $V$.  Let $M=r+k$.
For an increasing $M$-element list $J$, define
\begin{align}
 B_J(g)&=\operatorname{vol}
       (g^{-1}e_r,\ldots,g^{-1}e_1,g^{-1}z,e_J),
                                                     \label{rc-global-B}\\
 A_J(g)&=\det
       \begin{pmatrix}
        h(g\,\cdot,e_1)\\[-2pt]
        \vdots\\
        h(g\,\cdot,e_r)\\
        h(g\,\cdot,z)\\
        e_J^*
       \end{pmatrix}.
                                                     \label{rc-global-A}
\end{align}
The last $M$ rows in \eqref{rc-global-A}
are the coordinate covectors in the list $J$.
All matrix sizes are $2r+k+1$.
Set $\sigma=(-1)^{r(k+1)}$ and define
\begin{align}
 \mathcal R_Y F(g)
  &=-\sigma\sum_{i,J}(-1)^i b'_i B_J(g)\,
                 Q_{I_{0,i},J}(u_0)F(g),
                                                     \label{rc-global-R}\\
 \mathcal R_Y^- F(g)
  &=-\sum_{t,J}(-1)^t a'_t A_J(g)\,
                 Q_{J,J_{0,t}}(v-1)F(g).
                                                     \label{rc-global-Rminus}
\end{align}
These are the actual full-group coefficient
determinants in the source, with their constant
signs chosen explicitly.

\begin{lemma}\label{rc-cell-coefficients}
At $g=b(c)m$, all terms in
\eqref{rc-global-B} vanish except
$J=I_{0,t}$, and
\[
                 B_{I_{0,t}}(b(c)m)
                       =\sigma(-1)^t\overline{a(c)}\,a_t .
                                                     \tag{RC1}
\]
All terms in \eqref{rc-global-A} vanish except
$J=J_{0,s}$, and
\[
                 A_{J_{0,s}}(b(c)m)
                       =(-1)^s a(c)b_s .             \tag{RC2}
\]
\end{lemma}

\begin{proof}
In \eqref{rc-global-B} the first $r$ columns
are $\bar c_r e_r^\vee,\ldots,\bar c_1e_1^\vee$,
and the next column is $m^{-1}z$.
Any nonzero complementary coordinate determinant
must therefore contain all $U$-indices and all
but one of the $W$-indices.
Move the first $r$ columns past the other
$r+k+1$ columns, and then move the $m^{-1}z$
column past the $r$ coordinate columns of $U$.
These moves give the sign
$(-1)^{r(r+k+1)+r}=(-1)^{r(k+1)}$.
Expansion in its omitted $W$-coordinate gives
the extra sign $(-1)^t$.
This proves (RC1).

In \eqref{rc-global-A} the first $r$ rows are
$c_i e_i^*$, while the next row vanishes on
$U\oplus U^\vee$ and has $W$-coordinates $b_s$.
The complementary rows must therefore be the
dual indices and all but one $W$-index.
Expansion of that next row in column $r+s$
has sign $(-1)^{r+(r+s)}=(-1)^s$.
The remaining coordinate determinant is one.
This proves (RC2).
\end{proof}

\begin{theorem}\label{rc-actual-lifts}
The operators \eqref{rc-global-R} and
\eqref{rc-global-Rminus} have order at most
$r+k$ and polynomial coefficients of degree
$r+1$.  They are defined on the entire group.
They commute with left translation by $K_zN$
and right translation by $K_{z'}N$.
They and their Haar formal adjoints preserve
$\mathscr S_\varepsilon$.

For every $F$ satisfying \eqref{rt-whittaker},
their exact radial identities are
\begin{align}
 \gamma(c)(\mathcal R_YF)(b(c)m)
   &=\overline{a(c)}
       \bigl(t_{21}(Y)A_r(Y)+t_{22}(Y)B_r(Y)\bigr)\Phi,
                                                     \label{rc-radial-R}\\
 \gamma(c)(\mathcal R_Y^-F)(b(c)m)
   &=a(c)
       \bigl(t_{12}(Y)\overline A_r(Y)
                    +t_{22}(Y)\overline B_r(Y)\bigr)\Phi .
                                                     \label{rc-radial-Rminus}
\end{align}
The statements hold for every complex $Y$.
\end{theorem}

\begin{proof}
The coefficients $B_J$ have degree $r+1$
in $\overline g$, since $g^{-1}=Hg^*H$.
The coefficients $A_J$ have degree $r+1$
in $g$.  The order and polynomial assertions
follow from the sizes of the quantum minors.
The proof of Proposition~\ref{rt-polynomial}
also gives derivative bounds
$C_m(1+\|g\|)^{r+1}$ for these coefficients
and proves preservation of
$\mathscr S_\varepsilon$ by both operators
and their formal adjoints.

We give the full covariance argument to account
for a determinant factor which cannot simply
be omitted.
The vector
$\eta=e_r\wedge\cdots\wedge e_1\wedge z$
is fixed by $\bigwedge^{r+1}s_L$ for
$s_L\in K_zN$: the $U$-block is unipotent
and $s_Lz-z\in U$.
Hence $B_J(s_Lg)=B_J(g)$.
The corresponding wedge of covectors in
\eqref{rc-global-A} is likewise unchanged,
because unitarity replaces each fixed vector
by its image under $s_L^{-1}$.
Left translations commute with right
derivatives, proving left covariance.

Now let $s\in K_{z'}N$.
Put $P=U\oplus W$ and
$\ell(x)=h(x,z')$ for $x\in P$.
The action of $s$ preserves $P$ and satisfies
$\ell(sx)=\ell(x)$.
Its determinant on $P$ equals its determinant
on $V$: the actions on $U$ and on
$V/P$ have determinant one.
Thus the fixed row vector in the exterior
power,
\[
       \rho=\sum_i(-1)^i b'_i e_{I_{0,i}}
          =(-1)^r\iota_\ell
                  (e_1\wedge\cdots\wedge e_r
                         \wedge w_0\wedge\cdots\wedge w_k),
\]
satisfies
$\bigwedge^M(s^{-1})\rho=\det(s)^{-1}\rho$.
On the other side, the volume definition gives
\[
 \sum_J B_J(g)\det(s)_{J,L}
                  =\det(s) B_L(gs).
                                                     \tag{RC3}
\]
Indeed the left side is
$\operatorname{vol}(g^{-1}\eta\wedge s e_L)$;
apply $s^{-1}$ to every vector in this
volume.  Formula \eqref{rt-right-minor}
and (RC3) cancel the two determinant
factors and prove right covariance of
$\mathcal R_Y$.

For the second operator, the elementary
dual volume identity is
\[
 \sum_J\det(s^{-1})_{K,J}A_J(g)
                      =\det(s)^{-1}A_K(gs).
                                                     \tag{RC4}
\]
It follows either by applying $s^{-1}$ to
the coordinate rows in the full determinant,
or by the alternating Cauchy--Binet formula.
The fixed column covector is
\[
                  \kappa=\sum_t(-1)^t a'_t e_{J_{0,t}}^*.
 \]
On $V/U$, it is contraction of its standard
top-degree form by the vector $z'$.
The action of $s$ fixes this vector modulo
$U$ and has determinant $\det(s)$ on the
quotient.  Therefore
$\kappa\circ\bigwedge^M s=\det(s)\kappa$.
This factor cancels the one in (RC4) when
\eqref{rt-right-minor} is applied.
It proves right covariance of
$\mathcal R_Y^-$.

Finally (RC1), (RC2), and $\sigma^2=1$ give
\[
 \gamma(\mathcal R_YF)(b(c)m)
       =-\overline{a(c)}F_0,\qquad
 \gamma(\mathcal R_Y^-F)(b(c)m)
       =-a(c)F_0^- .
 \]
Corollary~\ref{rc-chain} proves
\eqref{rc-radial-R} and
\eqref{rc-radial-Rminus}.
The source sometimes writes the conjugate
transfer recurrence with its new coordinate
factor on the right.  Earlier-coordinate
operators commute with that factor, so it
is exactly the column used here.
\end{proof}

\paragraph{Scope.}
The four required polynomial differential
lifts and the regular spectral inverse have
now been constructed and matched with their
actual radial actions.  This is an algebraic
and differential statement.  The terminal
quantum relation used to cancel the two
mixed terms in the Fourier identity, and
the analytic passage to the Jacquet pairing
and its uniformly bounded generated class,
remain distinct obligations.  No integration
by parts for that pairing is inferred from
the ordinary Haar formal-adjoint calculation.


